% =====================================================================
% Resumo e Abstract
% =====================================================================
\clearpage
\phantomsection
\addcontentsline{toc}{section}{RESUMO}
\section*{RESUMO}

Sistemas de inteligência artificial vêm sendo incorporados a malhas de controle e a
cadeias de decisão em que uma saída fora de faixa deixa de ser um erro estatístico e
passa a ser uma falha de segurança. Este relatório apresenta os resultados do plano de
iniciação científica que investigou a aplicação de \textit{bounded model checking}
baseado em SMT, por meio do verificador ESBMC, a sete classes de artefatos de IA:
redes neurais quantizadas em ponto fixo, \textit{kernels} de inferência em C/C++,
blocos \textit{feed-forward} extraídos de modelos de linguagem (GPT-2 e LLaMA-7B),
um ciclo neuro-simbólico de geração e correção de código, um controlador PID sob ruído
não determinístico, uma política de aprendizado por reforço e, por fim, uma malha
fechada completa formada por planta \textit{cart-pole} e controlador DDPG quantizado.
A contribuição metodológica central não está no número de propriedades provadas, mas na
disciplina de evidência adotada: todo veredito relatado aqui é acompanhado da saída
bruta do \textit{solver}, do comando exato, das \textit{flags}, do limite de tempo e do
ambiente de execução, e \textit{timeout} é tratado como resultado \textbf{indeciso},
nunca como prova de segurança. Sob esse critério, dos doze \textit{harnesses}
reexecutáveis do repositório, nove foram provados, três produziram contraexemplo e
nenhum ficou indeciso. Duas propriedades da malha fechada foram diagnosticadas como
\textbf{vácuas} --- eram decididas sem que os 745 parâmetros do controlador
participassem da prova --- e foram reescritas de modo a depender efetivamente da rede.
A limitação central do método foi medida, e não estimada: provar $K$ passos da malha
fechada custa 0,3\,s em $K{=}1$ e 909\,s em $K{=}4$, de modo que a barreira prática do
BMC neste sistema está em quatro passos. Para contorná-la, o sistema acoplado foi
exportado como sistema de transição e submetido a \textit{model checking} indutivo
(IC3/PDR): no modelo sintético, o PDR produziu prova ilimitada em 0,37\,s onde o BMC
levou 909\,s para quatro passos; sobre o controlador real, porém, nenhum dos dois
métodos decide a instância. Conclui-se que a verificação formal de componentes de IA é
viável e informativa em escala modular, que a maior ameaça à validade desse tipo de
trabalho é a propriedade vácua, e que o próximo passo não é trocar de motor, mas mudar
a representação do problema para o nível de palavra.

\vspace{0.4cm}
\noindent\textbf{Palavras-chave}: verificação formal; ESBMC; \textit{bounded model
checking}; SMT; redes neurais quantizadas; IC3/PDR; segurança de sistemas de controle.

\clearpage
\phantomsection
\addcontentsline{toc}{section}{ABSTRACT}
\section*{ABSTRACT}
\begin{otherlanguage*}{english}

Artificial intelligence components are increasingly embedded in control loops and
decision pipelines where an out-of-range output is no longer a statistical error but a
safety failure. This report presents the results of an undergraduate research project
that investigated SMT-based bounded model checking, through the ESBMC verifier, applied
to seven classes of AI artifacts: fixed-point quantized neural networks, C/C++ inference
kernels, feed-forward blocks extracted from language models (GPT-2 and LLaMA-7B), a
neuro-symbolic code generation and repair loop, a PID controller under non-deterministic
sensor noise, a reinforcement learning policy, and a full closed loop comprising a
cart-pole plant and a quantized DDPG controller. The main methodological contribution is
not the number of properties proved but the evidence discipline adopted: every verdict
reported here is backed by the raw solver output, the exact command, the flags, the time
limit and the execution environment, and \textit{timeout} is treated as an
\textbf{undecided} outcome, never as proof of safety. Under this criterion, out of the
twelve reproducible harnesses in the repository, nine were proved, three produced
counterexamples and none remained undecided. Two closed-loop properties were diagnosed
as \textbf{vacuous} --- they were decided without the controller's 745 parameters taking
part in the proof --- and were rewritten so as to genuinely depend on the network. The
central limitation of the method was measured rather than estimated: proving $K$ steps of
the closed loop costs 0.3\,s at $K{=}1$ and 909\,s at $K{=}4$, so the practical BMC wall
for this system lies at four steps. To work around it, the coupled system was exported as
a transition system and submitted to inductive model checking (IC3/PDR): on the synthetic
model, PDR produced an unbounded proof in 0.37\,s where BMC took 909\,s for four steps;
on the real controller, however, neither method decides the instance. We conclude that
formal verification of AI components is feasible and informative at a modular scale, that
the greatest threat to validity in this line of work is the vacuous property, and that
the next step is not to switch engines but to change the problem representation to the
word level.

\vspace{0.4cm}
\noindent\textbf{Keywords}: formal verification; ESBMC; bounded model checking; SMT;
quantized neural networks; IC3/PDR; control system safety.

\end{otherlanguage*}
